\chapter{Estado del arte}

\section{Modelos de lenguaje con herramientas}

Los modelos de lenguaje modernos pueden resolver tareas complejas, pero su
utilidad aumenta cuando se les permite actuar sobre herramientas externas.
ReAct propone intercalar razonamiento y acciones para que el modelo consulte
fuentes externas, actualice planes y produzca trayectorias más interpretables
\cite{yao2023react}. Toolformer estudia la capacidad de los modelos para decidir
cuándo llamar APIs externas y cómo incorporar sus resultados \cite{schick2023toolformer}.

El sistema desarrollado adopta este patrón de forma controlada. Los agentes no
modifican libremente el entorno: reciben herramientas tipadas, como búsqueda web,
lectura y escritura de artefactos, recuperación de conocimiento o ejecución de
pruebas. El enrutador mantiene la responsabilidad de decidir qué etapa se
ejecuta y cuándo se acepta su resultado.

\section{Flujos multiagente}

Otra línea de trabajo relevante es la descomposición de tareas en agentes
especializados. AutoGen presenta un marco para construir aplicaciones con
múltiples agentes conversables, herramientas y participación humana
\cite{wu2023autogen}. ChatDev aplica la idea al desarrollo de software mediante
agentes especializados en diseño, codificación y prueba \cite{qian2024chatdev}.

El pipeline de esta memoria comparte la motivación de especialización, pero no
se basa en una conversación abierta entre agentes. La coordinación se realiza con
una máquina de estados explícita. Esta decisión reduce flexibilidad, pero mejora
trazabilidad y repetibilidad.

\begin{figure}[H]
\centering
\begin{verbatim}
Multiagente conversacional       Pipeline de esta memoria
--------------------------       ------------------------
Agente <-> Agente                Router -> Agente
Agente <-> Usuario               Revision humana
Control emergente                Estados explicitos
Historial conversacional         Artefactos versionados
\end{verbatim}
\caption{Diferencia entre conversación multiagente y pipeline orquestado.}
\end{figure}

\section{Recuperación aumentada}

Retrieval-Augmented Generation (RAG) combina memoria paramétrica del modelo con
memoria no paramétrica recuperada de una fuente externa \cite{lewis2020rag}.
Este enfoque es especialmente relevante cuando el sistema necesita conocimiento
actualizable, trazable o específico de un dominio.

La recuperación puede ser densa, basada en embeddings, o dispersa, basada en
coincidencia léxica. Dense Passage Retrieval muestra la utilidad de recuperar
pasajes mediante representaciones densas aprendidas \cite{karpukhin2020dpr}.
BM25 representa una base clásica y robusta para búsqueda textual probabilística
\cite{robertson2009bm25}. En el sistema implementado se usan ambas vías porque
resuelven problemas distintos: la búsqueda densa captura similitud semántica y
BM25 captura nombres, autores, símbolos, DOIs y términos exactos.

Cuando existen varios rankings, una técnica sencilla y efectiva es Reciprocal
Rank Fusion, que combina listas ordenadas sin necesidad de entrenamiento
\cite{cormack2009rrf}. Esta idea se usa para fusionar resultados densos,
lexicales y de grafo.

\section{RAG correctivo y reflexivo}

Un riesgo de RAG es inyectar contexto irrelevante o desactualizado. Self-RAG
plantea que la recuperación no debe ser indiscriminada y que el sistema debe
reflexionar sobre la utilidad de los pasajes recuperados \cite{asai2023selfrag}.
CRAG añade un evaluador de recuperación y acciones correctivas, incluyendo
búsqueda web cuando el corpus interno no contiene suficiente evidencia
\cite{yan2024crag}.

La arquitectura desarrollada adopta una variante práctica: cuando la confianza
de los resultados internos es alta, se evita trabajo extra; cuando es ambigua o
baja, se activa una clasificación correctiva y, si procede, una búsqueda web.

\section{Memoria a largo plazo y grafos}

Los agentes generativos de Park et al. almacenan experiencias, las sintetizan en
reflexiones y las recuperan dinámicamente para planificar comportamiento
\cite{park2023generativeagents}. MemoryBank propone memoria persistente con
actualización, importancia y olvido inspirado en la curva de Ebbinghaus
\cite{zhong2023memorybank}. MemGPT aborda el problema de la ventana de contexto
limitada mediante una memoria jerárquica gestionada fuera del prompt inmediato
\cite{packer2023memgpt}.

Para representar relaciones científicas, los grafos de conocimiento ofrecen un
modelo explícito de entidades y relaciones \cite{hogan2021knowledgegraphs}.
GraphRAG usa un grafo como índice sobre corpus privados para responder preguntas
globales \cite{edge2024graphrag}. HippoRAG combina LLMs, grafos de conocimiento
y Personalized PageRank para construir una memoria a largo plazo inspirada en el
hipocampo \cite{gutierrez2024hipporag}; Topic-sensitive PageRank es un
antecedente clásico de ranking sensible al contexto \cite{haveliwala2002topicsensitive}.

El sistema de esta memoria no replica exactamente GraphRAG ni HippoRAG. Sí adopta
su idea central: la memoria útil para agentes no debe ser solo una lista plana de
chunks, sino una combinación de hechos recuperables, entidades y relaciones.

\section{Generación de código y validación}

La generación automática de programas desde lenguaje natural se estudia en
trabajos como Codex y HumanEval \cite{chen2021codex} y Program Synthesis with
Large Language Models \cite{austin2021programsynthesis}. Estos trabajos muestran
que la corrección funcional debe evaluarse ejecutando el código, no solo leyendo
la respuesta del modelo.

SWE-agent refuerza la importancia de diseñar interfaces adecuadas para agentes
de software, incluyendo edición de archivos, navegación del repositorio y
ejecución de pruebas \cite{yang2024sweagent}. PICARD, aunque aplicado a
text-to-SQL, muestra el valor de restringir salidas para evitar estructuras
inválidas \cite{scholak2021picard}. Self-Refine y Reflexion aportan antecedentes
para refinar salidas mediante feedback y memoria verbal sin actualizar pesos
\cite{madaan2023selfrefine,shinn2023reflexion}.

En este trabajo, la validación se reparte en tres niveles: schemas y modelos
Pydantic para estructuras, revisión humana para decisiones científicas y tests
ejecutables para los modelos generados.
